\chapter{Content and Assets}
\label{ch:content-and-assets}

Real XNA's content model and CNA's are different, and --- unlike most of the gaps in this
book --- the story here has genuinely changed over the project's own history rather than
settling once. This chapter covers where it stands now: two real, actively-maintained content
formats, tried in a specific order, each with its own reasoning.

\section{Real XNA's own model, briefly}

Real XNA never ships raw asset files inside a game. An offline content-pipeline build step
compiles every texture, model, font, and effect into a proprietary binary format ---
\texttt{.xnb} --- which \cnaclass{ContentManager::Load<T>()} loads at runtime. FNA
reimplements a real \texttt{.xnb} reader, precisely so unmodified, historical XNA content
keeps working unchanged.

\section{CNA's own format: .cnj}
\label{sec:cnj}

CNA's primary, permanent strategic answer is not a \texttt{.xnb} reader at all, but its own
format: \texttt{.cnj} --- ``CNA content JSON,'' a lightweight, human-readable JSON sidecar
format sitting alongside a directly-loadable native asset file (a \texttt{.png}, a
\texttt{.wav}) or standing in for an asset with no native format of its own. This is a
permanent, deliberate strategic choice, not a stopgap awaiting \texttt{.xnb} support --- the
project's own naming history is worth knowing as a small, honest correction in its own right:
the format was originally called \texttt{.cnb} (``CNA content Binary/JSON''), a name that
misleadingly suggested a binary format mirroring XNA's own \texttt{.xnb}, despite always
having been plain JSON. It was renamed project-wide --- the file extension, the internal
envelope type, the version field, the registration template, and the format's own planning
document's task numbering all changed together --- specifically to stop implying something
the format never was.

A \texttt{.cnj} file is a small JSON envelope (a version field, a type name, and a reference
to a source file where applicable) validated against the requested C\texttt{++} type with a
strict version-equality policy, parsed by a real, from-scratch recursive-descent JSON parser.
\cnaclass{SpriteFont}, custom \cnaclass{Effect} parameters for the five stock effects
(Chapter~\ref{ch:stock-effects}), and \cnaclass{Model} (the fallback path from
\S\ref{sec:xnb-and-cnj-model} below) are all reachable through \texttt{.cnj} today, migrated
from older, bespoke loose-file extensions (\texttt{.font.json}, \texttt{.shader.json},
\texttt{.model.json}) that predated the unified format. A registration template lets a game
register several differently-named \texttt{.cnj} ``type'' strings that all resolve to the
same C\texttt{++} type, for data with no dedicated built-in reader at all --- the same
general-purpose extensibility mechanism Chapter~\ref{ch:shader-fx-gap} already noted CNA
deliberately chose over a scripted custom-reader mechanism, for the same reason: a plain,
type-safe C\texttt{++} registration API, not an embedded scripting language.

One asset type is deliberately \emph{not} part of the unified \texttt{.cnj} format, by an
explicit, reasoned decision: the Avatar system's \cnaclass{SkinnedModelEXT}
(Chapter~\ref{ch:avatar}) keeps its own separate \texttt{.skinnedmodel.json} extension, because
real cross-language tooling already depends on that specific extension name, because it already
carried the most mature test coverage of any loose-file reader in the project at the time
\texttt{.cnj} unification happened, and because it is explicitly a distinct,
non-\cnaclass{Model}-shaped system in the first place (Chapter~\ref{ch:avatar} covers exactly
why).

\subsection{Two real \texttt{.cnj} files, quoted directly}

Reading \texttt{cnj.md} (the format's own design reference) turns up two genuinely different
uses of the same envelope, worth seeing side by side rather than only described. First, a
fully self-contained document --- a real \cnaclass{SpriteFont} needs no native payload file at
all, since every field it needs (the glyph atlas reference, layout metrics, per-glyph bounds)
fits directly in JSON:

\begin{lstlisting}
{
  "cnjVersion": 1,
  "type": "SpriteFont",
  "texture": "Fonts/Consolas_atlas.png",
  "lineSpacing": 24,
  "spacing": 0.0,
  "defaultCharacter": "?",
  "glyphs": [
    { "char": 65, "source": [0, 0, 16, 24], "crop": [0, 0, 16, 24], "kerning": [1, 14, 1] }
  ]
}
\end{lstlisting}

Second, the \texttt{sourceFile} metadata-sidecar shape this chapter's resolution-order section
already named in the abstract, quoted here concretely:

\begin{lstlisting}
{
  "cnjVersion": 1,
  "type": "Texture2D",
  "sourceFile": "ahoj.png",
  "colorKey": [255, 0, 255],
  "premultipliedAlpha": true
}
\end{lstlisting}

The distinction between these two is exactly the \texttt{sourceFile} field's presence or
absence, and it changes what the rest of the document \emph{means}: without \texttt{sourceFile}
(the \cnaclass{SpriteFont} case), every other field is the actual data being described. With
\texttt{sourceFile} present (the \cnaclass{Texture2D} case), CNA loads \texttt{ahoj.png}
through the ordinary native-decode path first, and every other field in the \texttt{.cnj}
becomes import metadata layered on top of that already-decoded texture --- here, a magenta
\texttt{colorKey} (pixels matching \texttt{(255, 0, 255)} become transparent, the classic
sprite-sheet convention for games authored before per-pixel alpha channels were universal) and
an explicit premultiplied-alpha hint. This is genuinely new capability the format's own design
notes are explicit about, not a rename of something \cnaclass{Texture2D} could already do:
before \texttt{sourceFile} existed, a native-payload asset type could only ever load a
self-describing file exactly as authored, with no way to attach CNA-specific import intent
(color-key transparency, a premultiplied-alpha override) without baking it into the pixel data
itself or inventing an entirely separate file format just to carry one flag.

\section{The gap, reopened: a real, growing .xnb binary reader}
\label{sec:xnb-and-cnj-model}

Here is the correction this chapter most needs to make plainly: CNA's owner later decided
\texttt{.xnb} should \emph{also} become a real runtime format, alongside --- not instead of
--- \texttt{.cnj}. This is now a genuinely growing, real binary reader, not a planning
document: container parsing, binary primitives, and a first real \cnaclass{Texture2D} reader
are complete; real LZX decompression is implemented and verified against real compressed
fixtures, including a differential-testing pass against FNA's own reference decompressor
(run under Mono, the same technique Chapter~\ref{ch:what-is-cna} describes for
\texttt{tools/fna-reference/}) that found and fixed a real heap-buffer-overflow in the
process; and \cnaclass{SpriteFont}, all five stock effects, \cnaclass{SoundEffect} (at least
one real wave-format variant), \cnaclass{Song}, and the general
\texttt{ReadExternalReference<T>()} mechanism are all real, working \texttt{.xnb} readers
today. \cnaclass{Model}'s own \texttt{.xnb} reader --- described in
Chapter~\ref{ch:models-meshes} as already real and wire-compatible --- is exactly this
effort's own most complete result.

Writing or producing \texttt{.xnb} files remains permanently out of scope on both sides of
this story: CNA only ever needs to \emph{consume} \texttt{.xnb} files produced by real
XNA/MonoGame/FNA tooling, never to produce its own.

\section{The resolution order, precisely}

With both formats real, \cnaclass{ContentManager}'s own asset-resolution order matters, and
it is worth stating exactly: for a given asset name, \texttt{.xnb} is checked first; if no
\texttt{.xnb} file is found, the literal, caller-given path (a directly loadable native
asset, per \S\ref{sec:direct-loading} below) is tried next; only if neither of those resolves
does \texttt{.cnj} apply. This ordering is deliberate on both ends: \texttt{.xnb}, when
present, always wins outright, since it is the most capable and most binary-compatible
option; and \texttt{.cnj} is checked last specifically so it can act as an optional
\emph{metadata sidecar} for an otherwise-directly-loadable native asset (a \cnaclass{Texture2D}
with a \texttt{.cnj} file supplying a source-file override and a color key, for instance) as
well as a full replacement for an asset with no native format of its own --- not only a
mutually-exclusive alternative to loading the asset directly.

\section{Direct loading: the shape of Chapter~\ref{ch:first-game}'s smallest program}
\label{sec:direct-loading}

The one-line construction in Chapter~\ref{ch:first-game}'s minimal game ---
\texttt{Texture2D("assets/logo.png", getGraphicsDeviceProperty())} --- is the third, simplest
path this chapter describes: \texttt{"assets/logo.png"} is a literal, directly loadable file
path, resolved without any \texttt{.xnb} or \texttt{.cnj} involvement at all when neither is
present. This remains the simplest path for a small, from-scratch CNA project, and it is
exactly what the resolution order in the previous section falls through to only after the
other two have both been tried and found nothing.

\section{ContentManager: constructing it and calling Load}

The class every asset ultimately passes through, read directly from
\texttt{ContentManager.hpp}, has three constructors: \texttt{explicit
ContentManager(IServiceProvider*)}; a second form additionally taking
\texttt{const std::string\& rootDirectory}; and a \texttt{NOXNA}-tagged no-argument
\texttt{ContentManager()} that defaults \texttt{RootDirectory} to \texttt{"Content"} --- plus
the small set of methods a game actually calls day to day:

\begin{lstlisting}
// A Game already owns one ContentManager, built for it -- LoadContent() is where
// a real game reaches for it, per Chapter 5's ordering constraint:
ContentManager& content = getContentProperty();
content.setRootDirectoryProperty("Content");

auto logo = content.Load<Texture2D>("sprites/logo");   // no extension in the name
auto font = content.Load<SpriteFont>("fonts/hud");

content.Unload();  // releases every asset this manager has loaded
\end{lstlisting}

\texttt{Load<T>(const std::string\& assetName) -> T} is a template, not a virtual dispatch ---
the asset name carries no extension, since resolving one is exactly the three-way order this
chapter already covered (\texttt{.xnb}, then the literal path, then \texttt{.cnj}). Two
\texttt{NOXNA} extension points exist for cases the built-in readers don't cover:
\texttt{RegisterTypeReader<T>} (taking a
\texttt{std::unique\_ptr<LooseFileContentTypeReader<T>>}) for a whole new
loose-file type, and \texttt{RegisterCnjLoader<T>(name, CnjLoaderFn<T>)} for a new named
\texttt{.cnj} ``type'' string resolving to an existing C\texttt{++} type --- the mechanism
\S\ref{sec:cnj} already named as the reason a game never needs to invent its own ad-hoc JSON
format for a custom asset type.

\begin{sourcenote}
\textbf{A worked example: asking the content root what it would actually resolve.} Rather than
just describing the three-way resolution order in the abstract, CNA ships a real,
\texttt{NOXNA}-tagged introspection tool that answers it concretely for a whole content tree at
once: \texttt{ContentManager::GetContentManifest()} scans the content root (lazily, on first
call, via a private \texttt{RefreshContentManifest()}) and returns one
\cnaclass{ContentManifestEntry} per logical asset name, reading directly from
\texttt{ContentManifestEntry.hpp}:

\begin{lstlisting}
for (const ContentManifestEntry& entry : content.GetContentManifest()) {
    std::cout << entry.relativePath << ": ";
    if (entry.hasXnb) {
        std::cout << "resolves via .xnb"; // wins outright when present
    } else if (!entry.nativeExtensions.empty()) {
        std::cout << "resolves via direct load (" << entry.nativeExtensions[0] << ")";
    } else if (entry.hasCnj) {
        std::cout << "resolves via .cnj only";
    }
    std::cout << "\n";
}
\end{lstlisting}

Each \cnaclass{ContentManifestEntry} exposes exactly the three booleans/lists this chapter's
resolution-order section depends on --- \texttt{hasXnb}, \texttt{hasCnj}, and
\texttt{nativeExtensions} --- plus, when a \texttt{.xnb} file is present and readable,
\texttt{xnbReaderNames}: the canonical reader-type names its own type-reader table references,
letting a companion method, \texttt{GetXnbReaderUsageSummary()}, aggregate across an entire
content tree which reader types are actually exercised and whether
\cnaclass{ContentTypeReaderManager} currently has a registered reader for each one --- a real,
runnable way to answer ``which of my compiled XNA content's readers does CNA not support yet''
without reading a single byte of any individual \texttt{.xnb} file by hand.
\end{sourcenote}

\section{What this means for porting}

Chapter~\ref{ch:migration-guide}'s migration guide covers this in full practical detail, but the
short version, now that this chapter has given the corrected picture, is worth stating here:
porting an existing XNA/FNA game's compiled \texttt{.xnb} content is, for the specific asset
types \S\ref{sec:xnb-and-cnj-model} named as complete (textures, sprite fonts, stock effects,
sound effects, songs, and models), increasingly a matter of pointing CNA at the original,
unmodified compiled assets and letting the real binary reader do the work --- not, as an
earlier version of this book's own understanding assumed, a mandatory translation into
hand-authored JSON. For asset types the \texttt{.xnb} reader has not reached yet, or a custom
compiled \texttt{.fx} effect (permanently out of scope for the reason
Chapter~\ref{ch:shader-fx-gap} already gave), \texttt{.cnj} or a direct native-format
conversion remains the path.

\section{Where this leaves the rest of Part~III}

Every class in the chapters ahead --- \cnaclass{Texture2D}, \cnaclass{Model}, the stock
effects --- is describable purely in terms of its runtime API, independent of which of this
chapter's three loading paths actually produced a live instance of it. This chapter's
distinction matters specifically at the boundary between ``I have an asset file'' and ``I have
a live CNA object,'' which is exactly the boundary Chapter~\ref{ch:migration-guide}'s migration guide and the
\texttt{cna-samples} / \texttt{cna-examples} chapter (Chapter~\ref{ch:samples-and-examples})
come back to when the asset in
question was originally authored for real XNA rather than for CNA from scratch.
